\def\ProblemCode{GEOMETRY}
\def\ProblemTitle{Hình học cơ bản}
\makeproblem{\ProblemCode}{\ProblemTitle}

Cho $n$ điểm $P_1,P_2,\dots,P_n$ trên mặt phẳng tọa độ, được đánh số từ $1$ đến $n$.  
Hãy xử lý $q$ truy vấn hình học trên các điểm này.

Các truy vấn có một trong các dạng sau:

\begin{itemize}
    \item \boxed{\texttt{ADD}\ i\ j}: In ra tọa độ của $\vec{P_i} + \vec{P_j}$.

    \item \boxed{\texttt{SCALE}\ i\ k}: In ra tọa độ của $k \cdot \vec{P_i}$.

    \item \boxed{\texttt{EQUAL}\ i\ j}: Kiểm tra hai điểm $P_i$ và $P_j$ có trùng nhau hay không.

    \item \boxed{\texttt{DOT}\ i\ j}: In ra tích vô hướng $\vec{P_i} \cdot \vec{P_j}$.

    \item \boxed{\texttt{CROSS}\ i\ j}: In ra tích có hướng $\vec{P_i} \times \vec{P_j}$.

    \item \boxed{\texttt{CCW}\ a\ b\ c}: Xác định hướng quay của ba điểm $P_a, P_b, P_c$.
    
    In ra:
    \begin{itemize}
        \item \boxed{\texttt{1}} nếu $P_a \rightarrow P_b \rightarrow P_c$ quay ngược chiều kim đồng hồ.
        \item \boxed{\texttt{-1}} nếu quay theo chiều kim đồng hồ.
        \item \boxed{\texttt{0}} nếu ba điểm thẳng hàng.
    \end{itemize}

    \item \boxed{\texttt{DIST}\ i\ j}: In ra khoảng cách giữa $P_i$ và $P_j$.

    \item \boxed{\texttt{ONSEG}\ p\ a\ b}: Kiểm tra $P_p$ có nằm trên đoạn $P_aP_b$ hay không $(P_a \ne P_b)$.

    \item \boxed{\texttt{DISTSEG}\ p\ a\ b}: In ra khoảng cách từ $P_p$ đến đoạn $P_aP_b$. $(P_a \ne P_b)$.

    \item \boxed{\texttt{PROJECT}\ p\ a\ b}: In ra hình chiếu vuông góc của $P_p$ lên đường thẳng $P_aP_b$ $(P_a \ne P_b)$.

    \item \boxed{\texttt{REFLECT}\ p\ a\ b}: In ra điểm đối xứng của $P_p$ qua đường thẳng $P_aP_b$ $(P_a \ne P_b)$.

    \item \boxed{\texttt{INTERSECT}\ a\ b\ c\ d}: Kiểm tra hai đoạn $P_aP_b$ và $P_cP_d$ có cắt nhau hay không.
    
    Nếu chúng cắt nhau tại đúng một điểm $P$, in \texttt{YES} rồi tọa độ $P$.
    Nếu không, in \texttt{NO}.
    
    \textbf{Lưu ý:} Nếu hai đoạn thẳng trùng nhau hoặc chồng lấp trên vô số điểm,
    truy vấn này cũng được xem là \texttt{NO}.
\end{itemize}

% ===== INPUT DESCRIPTION =====
\inputDesc{}

\begin{itemize}
    \item Dòng đầu tiên gồm hai số nguyên $n, q\ (1 \le n, q \le 10^5)$.
    \item $n$ dòng tiếp theo, dòng thứ $i$ chứa hai số thực $x_i, y_i\ (-10^6 \le x_i,y_i \le 10^6)$ là tọa độ của điểm $i$. Phần thập phân của mỗi tọa độ có nhiều nhất $6$ chữ số.
    \item $q$ dòng tiếp theo, mỗi dòng chứa một truy vấn theo các dạng đã mô tả ở trên.
\end{itemize}

% ===== OUTPUT DESCRIPTION =====
\outputDesc{}

\begin{itemize}
    \item Với mỗi truy vấn, in ra kết quả tương ứng trên một dòng.
\end{itemize}

Các giá trị thực sẽ được chấp nhận nếu sai số tuyệt đối không vượt quá $10^{-3}$,
tức là nếu đáp án đúng là $x$ và bạn in ra $y$ thì $|x-y| \le 10^{-3}$.

\vspace{0.5em}

\ExampleBox{1}{
5 15\\
0 0\\
2 0\\
2 2\\
0 2\\
1 3\\
ADD 2 4\\
SCALE 3 1.5\\
EQUAL 5 3\\
DOT 2 4\\
CROSS 2 4\\
CCW 3 4 5\\
DIST 1 3\\
ONSEG 5 3 4\\
DISTSEG 5 1 4\\
DISTSEG 3 1 4\\
PROJECT 5 1 4\\
REFLECT 5 1 4\\
REFLECT 5 1 3\\
INTERSECT 1 3 2 4\\
INTERSECT 4 5 1 2
}{
2.000000000000 2.000000000000\\
3.000000000000 3.000000000000\\
NO\\
0.000000000000\\
4.000000000000\\
-1\\
2.828427124746\\
NO\\
1.414213562373\\
2.000000000000\\
0.000000000000 3.000000000000\\
-1.000000000000 3.000000000000\\
3.000000000000 1.000000000000\\
YES 1.000000000000 1.000000000000\\
NO
}{}